% GENERATED FILE. Edit canonical lessons, then run npm run generate:books.
% canonical-source-hash: fnv1a64:2f081598f3090c5c
% canonical-lessons: LA-C120-saccus, LA-C120-cista, LA-C120-pulvinus, LA-C120-stragulum, LA-C120-mappa

\chapter{Sack, Box, Cushion, Blanket, Napkin}
\label{ch:la-sack-box-cushion-blanket-napkin}
\hlchaptermodality{120}

\begin{chapteropening}
\textbf{By the end of this chapter:} \emph{I can say a sack, a box, a cushion, a blanket and a napkin.}

Complete the last lesson of chapter 120: I can say a sack, a box, a cushion, a blanket and a napkin.
\end{chapteropening}

\section[saccus, saccī]{saccus, saccī — a sack, a bag}

\label{lesson:LA-C120-saccus}



\begin{quote}
Before the new one: say the Latin for next to, close by, then the Latin for opposite, against.
\end{quote}

\subsection*{You'll want to know: saccus, saccī}
\textbf{saccus, saccī} — \textquotedblleft{}a sack, a bag\textquotedblright{}.

\begin{cousinweb}
\textbf{saccus} was borrowed from Greek \emph{sakkos}, \textquotedblleft{}a bag of coarse cloth\textquotedblright{} (Lewis \& Short), which Greek took from a Semitic language. English \emph{sack} comes from it.
\end{cousinweb}

\begin{grammarlens}[title={What you've built}]
One more word for this chapter.
\end{grammarlens}

\subsection*{Your turn}
\begin{itemize}
\raggedright
  \item \emph{Say it:} \emph{saccus}
  \item \emph{Say it:} \emph{saccus}, once more
  \item \emph{Say it:} say \emph{saccus}
  \item \emph{Recall:} say the Latin for next to, close by, then the Latin for opposite, against, then say \emph{saccus} again
\end{itemize}

\begin{tcolorbox}[breakable,colback=teal!4,colframe=teal!35!black,title={Before you move on}]
What does saccus, saccī mean? (\textquotedblleft{}A sack, a bag\textquotedblright{}.) Say it once more.
\end{tcolorbox}
\section[cista, cistae]{cista, cistae — a box, a chest}

\label{lesson:LA-C120-cista}



\begin{quote}
Before the new one: say the Latin for opposite, against, then the Latin for a sack.
\end{quote}

\subsection*{You'll want to know: cista, cistae}
\textbf{cista, cistae} — \textquotedblleft{}a box, a chest\textquotedblright{}.

\begin{cousinweb}
\textbf{cista} was borrowed from Greek \emph{kistē}, \textquotedblleft{}a box, a basket\textquotedblright{} (Lewis \& Short). English \emph{chest} was borrowed from it very early.
\end{cousinweb}

\begin{grammarlens}[title={What you've built}]
One more word for this chapter.
\end{grammarlens}

\subsection*{Your turn}
\begin{itemize}
\raggedright
  \item \emph{Say it:} \emph{cista}
  \item \emph{Say it:} \emph{cista}, once more
  \item \emph{Say it:} say \emph{cista}
  \item \emph{Recall:} say the Latin for opposite, against, then the Latin for a sack, then say \emph{cista} again
\end{itemize}

\begin{tcolorbox}[breakable,colback=teal!4,colframe=teal!35!black,title={Before you move on}]
What does cista, cistae mean? (\textquotedblleft{}A box, a chest\textquotedblright{}.) Say it once more.
\end{tcolorbox}
\section[pulvīnus, pulvīnī]{pulvīnus, pulvīnī — a cushion, a pillow}

\label{lesson:LA-C120-pulvinus}



\begin{quote}
Before the new one: say the Latin for a sack, then the Latin for a box.
\end{quote}

\subsection*{You'll want to know: pulvīnus, pulvīnī}
\textbf{pulvīnus, pulvīnī} — \textquotedblleft{}a cushion, a pillow\textquotedblright{}.

\begin{cousinweb}
\textbf{pulvīnus} has no sure origin; the link with \emph{pulvis}, \textquotedblleft{}dust\textquotedblright{}, that EtymDB 2.1 records is doubtful.
\end{cousinweb}

\begin{grammarlens}[title={What you've built}]
One more word for this chapter.
\end{grammarlens}

\subsection*{Your turn}
\begin{itemize}
\raggedright
  \item \emph{Say it:} \emph{pulvīnus}
  \item \emph{Say it:} \emph{pulvīnus}, once more
  \item \emph{Say it:} say \emph{pulvīnus}
  \item \emph{Recall:} say the Latin for a sack, then the Latin for a box, then say \emph{pulvīnus} again
\end{itemize}

\begin{tcolorbox}[breakable,colback=teal!4,colframe=teal!35!black,title={Before you move on}]
What does pulvīnus, pulvīnī mean? (\textquotedblleft{}A cushion, a pillow\textquotedblright{}.) Say it once more.
\end{tcolorbox}
\section[strāgulum, strāgulī]{strāgulum, strāgulī — a blanket, a bedspread}

\label{lesson:LA-C120-stragulum}



\begin{quote}
Before the new one: say the Latin for a box, then the Latin for a cushion.
\end{quote}

\subsection*{You'll want to know: strāgulum, strāgulī}
\textbf{strāgulum, strāgulī} — \textquotedblleft{}a blanket, a bedspread\textquotedblright{}.

\begin{cousinweb}
\textbf{strāgulum} is the neuter of \emph{strāgulus}, \textquotedblleft{}for spreading out\textquotedblright{} (Lewis \& Short), built on the root of \emph{sternō}, \textquotedblleft{}to spread, to lay out\textquotedblright{}.
\end{cousinweb}

\begin{grammarlens}[title={What you've built}]
One more word for this chapter.
\end{grammarlens}

\subsection*{Your turn}
\begin{itemize}
\raggedright
  \item \emph{Say it:} \emph{strāgulum}
  \item \emph{Say it:} \emph{strāgulum}, once more
  \item \emph{Say it:} say \emph{strāgulum}
  \item \emph{Recall:} say the Latin for a box, then the Latin for a cushion, then say \emph{strāgulum} again
\end{itemize}

\begin{tcolorbox}[breakable,colback=teal!4,colframe=teal!35!black,title={Before you move on}]
What does strāgulum, strāgulī mean? (\textquotedblleft{}A blanket, a bedspread\textquotedblright{}.) Say it once more.
\end{tcolorbox}
\section[mappa, mappae]{mappa, mappae — a napkin, a cloth}

\label{lesson:LA-C120-mappa}



\begin{quote}
Before the new one: say the Latin for a cushion, then the Latin for a blanket.
\end{quote}

\subsection*{You'll want to know: mappa, mappae}
\textbf{mappa, mappae} — \textquotedblleft{}a napkin, a cloth\textquotedblright{}.

\begin{cousinweb}
\textbf{mappa} came from Punic, the language of Carthage: Quintilian already called \emph{mappa} a Punic word (Lewis \& Short). English \emph{map} comes from it, through \emph{mappa mundi}, \textquotedblleft{}cloth of the world\textquotedblright{}.
\end{cousinweb}

\begin{grammarlens}[title={What you've built}]
One more word for this chapter.
\end{grammarlens}

\subsection*{Your turn}
\begin{itemize}
\raggedright
  \item \emph{Say it:} \emph{mappa}
  \item \emph{Say it:} \emph{mappa}, once more
  \item \emph{Say it:} say \emph{mappa}
  \item \emph{Recall:} say the Latin for a cushion, then the Latin for a blanket, then say \emph{mappa} again
\end{itemize}

\begin{tcolorbox}[breakable,colback=teal!4,colframe=teal!35!black,title={Before you move on}]
What does mappa, mappae mean? (\textquotedblleft{}A napkin, a cloth\textquotedblright{}.) Say it once more.
\end{tcolorbox}
